% فهرست اختصارات
\chapter*{فهرست اختصارات}
\phantomsection
\addcontentsline{toc}{chapter}{فهرست اختصارات}

در این پایان‌نامه به‌منظور سهولت نگارش و اختصار متن، از کوته‌نوشت‌های استاندارد زیر استفاده شده است:

\begin{center}
\begin{tabular}{|c|c|}
\hline
\textbf{علامت اختصاری} & \textbf{شرح فارسی} \\
\hline
\lr{ANN} & شبکه عصبی مصنوعی \\
\hline
\lr{BC} & شرط مرزی \\
\hline
\lr{CPU} & واحد پردازش مرکزی \\
\hline
\lr{GPU} & واحد پردازش گرافیکی \\
\hline
\lr{IC} & شرط اولیه \\
\hline
\lr{L-BFGS} & روش شبه‌نیوتنی با حافظه محدود \\
\hline
\lr{MLP} & پرسپترون چندلایه \\
\hline
\lr{MSE} & میانگین مربعات خطا \\
\hline
\lr{ODE} & معادله دیفرانسیل معمولی \\
\hline
\lr{PDE} & معادله دیفرانسیل جزئی \\
\hline
\lr{PINN} & شبکه عصبی آگاه از فیزیک \\
\hline
\lr{ReLU} & واحد خطی یکسوشده \\
\hline
\lr{RK} & رانگه-کوتا \\
\hline
\lr{SSE} & مجموع مربعات خطا \\
\hline
\end{tabular}
\end{center}
